\section{Methods}
\label{sec:methods}

\subsection{The $2 \times 2$ Epistemic Uncertainty Formulation}
We formalize conversational RAG interaction as a sequential decision process over a state $s = (\mathbf{x}, \mathcal{D}, \mathcal{Q}, \mathcal{R})$, where $\mathbf{x} \in \mathbb{R}^{d_q}$ represents the user's initial query, $\mathcal{D} = \{\bm{d}_1, \dots, \bm{d}_K\}$ is an unordered multiset of $K$ retrieved passages, $\mathcal{Q} = \{\bm{q}_1, \dots, \bm{q}_Q\}$ is a pool of candidate clarification questions specific to each $s$, and $\mathcal{R} = \{\bm{r}_1, \dots, \bm{r}_L\}$ is a pool of candidate search query rewrites.

We define two binary variables that govern the state:
\begin{align*}
    P &\in \{0, 1\} \quad \text{(Prompt Certainty: } 1 = \text{unambiguous}, 0 = \text{ambiguous)} \\
    E &\in \{0, 1\} \quad \text{(Evidence Sufficiency: } 1 = \text{verified/sufficient}, 0 = \text{missing/contradictory)}
\end{align*}
The Cartesian product $P \times E$ forms four quadrants, each mapped directly to an optimal macro action $a^* \in \mathcal{A} = \{\texttt{ANSWER}, \texttt{ASK}, \texttt{SEARCH}, \texttt{ABSTAIN}\}$:
\begin{itemize}
    \item \textbf{Quadrant 0 ($P=1, E=1 \to \texttt{ANSWER}$):} The user query is specific, and retrieved passages contain consistent, factual proof. The system generates a direct response immediately.
    \item \textbf{Quadrant 1 ($P=0, E=1 \to \texttt{ASK}$):} The query admits multiple valid interpretations, and 5retrieved passages cover them all. The optimal action is to ask a clarification question $q^* \in \mathcal{Q}$ that discriminates between competing interpretations.
    \item \textbf{Quadrant 2 ($P=1, E=0 \to \texttt{SEARCH}$):} The query is clear, but retrieved passages are contradictory or missing bridge hops. Clarifying with the user is futile; the system must reformulate search query $r^* \in \mathcal{R}$ to retrieve missing evidence.
    \item \textbf{Quadrant 3 ($P=0, E=0 \to \texttt{ABSTAIN}$):} The query is underspecified, and the underlying corpus contains unresolvable gaps. The system abstains to avoid a failure on both fronts.
\end{itemize}

\subsection{Counterfactual Value of Information (VOI) Framework}
Let $R(s) \in [0, 1]$ represent the immediate decision regret (the expected error penalty) if the system attempts to answer in state $s$. If the system executes action $a$, the expected post-action regret is $\mathbb{E}[R(s \mid a)]$. The expected regret reduction is defined as $\Delta R(s,a) = \max(0, R(s) - \mathbb{E}[R(s \mid a)])$.

In decision theory\cite{howard1966information}, actions incur operational costs. We model two distinct penalties:
\begin{itemize}
    \item \textbf{User Inconvenience Cost ($\lambda = 0.040$):} Asking the user a clarification question incurs cognitive overhead and conversational friction.
    \item \textbf{Computational Retrieval Cost ($\gamma = 0.020$):} Executing a secondary search query incurs compute latency and index retrieval cost.
\end{itemize}
The net Value of Information for candidate clarification question $q \in \mathcal{Q}$ and search rewrite $r \in \mathcal{R}$ is:
\begin{align*}
    \operatorname{VOI}(q) &= \Delta R(s,q) - \lambda \cdot \operatorname{Cost}(q) \\
    \operatorname{VOI}(r) &= \Delta R(s,r) - \gamma \cdot \operatorname{Cost}(r)
\end{align*}
Here, $\operatorname{Cost}(x)\in (0,1]$ models the cognitive effort required by the user (for $x\in\mathcal{Q}$) or the computational cost of retrieval (for $x\in\mathcal{R}$).

In Quadrant 1, appropriate and strong clarification questions yield $\Delta R(s,q) \approx 0.25 - 0.35$, yielding a positive net $\operatorname{VOI}(q) > 0$. In contrast, surface distractor questions (e.g., broad formatting questions) have zero regret reduction ($\Delta R = 0$), resulting in strictly negative net $\operatorname{VOI}(q) = -\lambda \cdot \operatorname{Cost}(q) < 0$. In non-ASK states (Quadrants 0, 2, 3), $\Delta R(q) = 0$ for all questions, ensuring that clarification questions are never selected over answering or searching.

\subsection{Neural Policy Architectures}
We compare six different model architectures to evaluate how different deep learning components impact performance:

\subsubsection{Model A: Retrieval Heuristic Baseline}
Model A reflects how simple production RAG systems make decisions using basic rules. It computes action scores using linear thresholds on top passage similarity and query length:
\begin{equation}
    \text{Score}_{\texttt{ANSWER}} = (\max_k \operatorname{sim}(\bm{x}, \bm{d}_k) - 0.65) \times 0.5, \quad \text{Score}_{\texttt{ASK}} = (0.50 - \operatorname{len}(\bm{x})) \times 0.4 \tag{A}
\end{equation}
Because Model A only looks at passages individually, it cannot detect when two retrieved passages contradict each other or when a query needs domain-specific clarification.

\subsubsection{Model B: NLI-Based LLM Emulator}
To represent how an instruction-tuned LLM behaves when prompted zero-shot, Model B uses Natural Language Inference (NLI) signals instead of simple vector similarity \cite{kuhn2023semantic, manakul2023selfcheckgpt, honovich2022true}. It estimates prompt certainty from entity specificity and multi-intent likelihood, and evidence certainty by checking whether passages entail the answer while penalizing contradictions:
\begin{align*}
    P_{\text{cert}} &= \operatorname{clip}(0.65 \cdot \text{EntitySpec} + 0.35(1 - \text{MultiIntent}), 0.05, 0.95) \\
    E_{\text{cert}} &= \operatorname{clip}(\text{Entail} - 0.35 \cdot \text{Contra}, 0.05, 0.95) \tag{B}
\end{align*}
Following empirical findings on LLM calibration \cite{xiong2024can}, Model B also incorporates verbalized overconfidence (outputting confidence scores biased towards $70\% - 95\%$) and a tendency to prefer polite, formatting-based clarification questions over deeper semantic disambiguation.

\subsubsection{Model C: Pooled MLP Baseline (Ablation 1)}
Model C is a standard feed-forward network. It averages (mean-pools) the feature vectors across all $K=6$ retrieved passages into a single vector, and passes this through a 4-layer MLP (hidden dimension 128). This serves as our baseline to test whether an attention mechanism across passages is actually needed.

\subsubsection{Model D: Cross-Passage Set Transformer}
Retrieved passages do not have an inherent sequence order. Model D processes them as an unordered set using a Set Transformer \cite{lee2019set}. It uses Multi-Head Attention Blocks (MAB) to let the query attend to all passages:
\begin{equation}
    \bm{H} = \operatorname{MAB}(\bm{X}, \bm{D}) = \operatorname{LayerNorm}(\bm{X} + \operatorname{MultiHead}(\bm{X}, \bm{D}, \bm{D})) \tag{D}
\end{equation}
followed by multi-head pooling over passages. Crucially, this allows the network to compare passages directly against one another, making it easy to identify when two retrieved passages contain conflicting facts.

\subsubsection{Model E: Latent Intent Set Transformer}
When a user asks an ambiguous question, there is usually a specific constraint or interpretation they had in mind that was left unsaid (such as a specific year, movie version, or metric). Model E builds on Model D by adding a head that predicts a Gaussian distribution (mean $\bm{\mu}$ and variance $\bm{\sigma}^2$) over these possible underlying user intents:
\begin{equation*}
    q(\bm{\theta} \mid s) = \mathcal{N}(\bm{\mu}(s), \operatorname{diag}(\bm{\sigma}^2(s)))
\end{equation*}
trained with Negative Log-Likelihood (NLL). When a query is underspecified, the predicted variance is naturally larger, signaling to the question-scoring head that it should choose a clarification question that targets this uncertainty.

\subsubsection{Model F: Full Proposed Dual-Uncertainty Architecture}
Model F brings together the Set Transformer, the latent intent head, and two additional safeguards for reliable deployment:
\begin{enumerate}
    \item \textbf{Deep Ensembles ($K_{\text{ens}}=3$):} Using three independent prediction heads \cite{lakshminarayanan2017simple} to measure model disagreement and prevent overconfident mistakes.
    \item \textbf{Temperature Scaling:} Calibrating output probabilities on the validation set \cite{guo2017calibration} so that the model's reported confidence accurately reflects its real-world accuracy.
\end{enumerate}

\subsection{Multi-Task Loss and Distractor Suppression}
The neural models are trained end-to-end with a multi-task loss that teaches the network both high-level routing and fine-grained action ranking:
\begin{equation*}
    \mathcal{L}_{\text{total}} = 2.5 \mathcal{L}_{\text{macro\_reg}} + 0.5 \mathcal{L}_{\text{macro\_ce}} + 2.5 \mathcal{L}_{\text{q\_reg}} + 2.0 \mathcal{L}_{\text{q\_ce}} + 2.0 \mathcal{L}_{\text{q\_margin}} + 1.5 \mathcal{L}_{\text{r\_reg}} + 0.5 \mathcal{L}_{\text{quad}} + 1.0 \mathcal{L}_{\text{regret}}
\end{equation*}
Here, $\mathcal{L}_{\text{macro}}$ supervises the choice between $\{\texttt{ANSWER}, \texttt{ASK}, \texttt{SEARCH}, \texttt{ABSTAIN}\}$, while $\mathcal{L}_{\text{q\_reg}}$ and $\mathcal{L}_{\text{r\_reg}}$ train the heads to predict the expected VOI of candidate questions and search rewrites.

\subsubsection{Fixing Distractor Question Bias}
In preliminary experiments, we observed a subtle recurring failure: surface distractor questions (like \textit{"Would you like a short summary or a long list?"}) have a lower cost than actual clarification questions. Because of this, models often assigned higher scores to distractors when the query was already clear. We solved this with two specific training constraints:
\begin{enumerate}
    \item \textbf{Selective Cross-Entropy on Questions:} We only apply classification cross-entropy ($\mathcal{L}_{\text{q\_ce}}$) on states where clarification is genuinely needed (Quadrant 1), preventing the model from learning to ask distractors during normal answer or search turns.
    \item \textbf{Margin-Based Distractor Penalty:} We enforce a ranking margin requiring real semantic questions to score at least 0.20 higher than distractor questions:
    \begin{equation*}
        \mathcal{L}_{\text{q\_margin}} = \frac{1}{B} \sum_{i=1}^B \sum_{j \in \text{distractors}} \max\left(0, s(q_{i, j}) + 0.20 - \max_{m \in \text{non-distractor}} s(q_{i, m})\right)
    \end{equation*}
    This penalty guarantees that even when an ambiguous query arises, the model strictly prefers meaningful clarification over trivial formatting questions.
\end{enumerate}